%=====================================================================
\chapter{Answers to the Checkpoints}\label{app:answers}
%=====================================================================
\normalfigures

Three answers per chapter, in order. Attempt the checkpoint from memory before
reading these, and argue with them where you disagree --- several are judgement
calls rather than facts, and they say so.

% The slide pipeline parses this appendix: \ansch{<number> --- <title>} must be
% followed immediately by an ans environment of exactly three items.
\newcommand{\ansch}[1]{\subsection*{\color{primary}Chapter #1}}
\newenvironment{ans}
  {\begin{enumerate}[leftmargin=*, itemsep=4pt, topsep=3pt,
                     label=\textbf{\arabic*.}]}
  {\end{enumerate}}

%---------------------------------------------------------------------
\ansch{1 --- The Role of Algorithms}
\begin{ans}
  \item $n$ grew by 4 and the time by 16, so the time grew as $n^{2}$:
        $\bigTh{n^{2}}$. The reasoning is the ratio test of Section~1.4 --- for
        $T = cn^{2}$, $T(4n)/T(n) = 16$ independent of $c$, so the constant
        never has to be known. One check is weak evidence; three doublings in
        a row would be convincing.
  \item A: $(10^{4})^{3}/10^{9} = 10^{12}/10^{9} = \mathbf{1000}$ s. B:
        $100 \times (10^{4})^{2}/10^{7} = 10^{10}/10^{7} = \mathbf{1000}$ s.
        They \textbf{tie exactly} at $n = 10^{4}$. Below it B wins, above it A
        loses --- the crossover is at that point, which is the real content of
        the question.
  \item \textbf{Harmless:} comparing insertion sort with merge sort at
        $n = 32{,}000$, where the orders of growth differ by a factor of 50 and
        a cache factor of a few cannot reverse the ranking.
        \textbf{Not harmless:} comparing two $\bigTh{n\lg n}$ sorts, where the
        orders of growth are equal and memory behaviour is the whole
        difference --- \chref{mergesort} measured \texttt{std::sort} beating
        merge sort by 39\% for exactly this reason.
\end{ans}

%---------------------------------------------------------------------
\ansch{2 --- Input Size and the Three Cases}
\begin{ans}
  \item Yes. The input is $n^{2}$ numbers, so writing $N = n^{2}$ for the true
        input size gives $\bigTh{n^{3}} = \bigTh{N^{1.5}}$ --- polynomial in
        $N$. Indeed $\bigTh{n^{3}}$ is only $n$ times the time needed to
        \emph{read} the input, which is why matrix multiplication is considered
        nearly efficient rather than slow.
  \item Because $(n+1)/2$ and $n$ differ by a constant factor of 2, and
        $\bigTh{\cdot}$ discards constants. Both are $\bigTh{n}$. The average
        case differs from the worst only in the constant, exactly as insertion
        sort's does (Section~2.2), and that is the common situation.
  \item \textbf{``Averaged over what distribution of inputs?''} Without it the
        claim is unfalsifiable. The usual answer --- a uniform random
        permutation --- is very often false of real data, which arrives partly
        sorted; and if the answer is instead ``expected over the algorithm's
        own randomness'' then the bound is much stronger and holds for your
        data too.
\end{ans}

%---------------------------------------------------------------------
\ansch{3 --- Asymptotic Notation}
\begin{ans}
  \item For $n \ge 1$ we have $n^{2} \le n^{3}$ and $1 \le n^{3}$, so
        $5n^{3} + 200n^{2} + 17 \le (5+200+17)n^{3} = 222n^{3}$. Take
        $c = 222$, $n_{0} = 1$. At $n = 1$: LHS $= 222$, RHS $= 222$, so the
        inequality holds with equality. A sharper pair such as $c = 6$,
        $n_{0} = 201$ is equally valid --- the definition asks only that
        \emph{some} pair exist.
  \item \textbf{Yes.} Given any $c > 0$ we need $n^{2} < c\,n^{2}\lg n$ for all
        large $n$, i.e.\ $1 < c\lg n$, i.e.\ $n > 2^{1/c}$. So
        $n_{0} = \lceil 2^{1/c}\rceil + 1$ works for every $c$, which is the
        definition of $o$. (Note $2^{1/c}$ is large when $c$ is small, as it
        must be --- but it is finite, and that is all that is required.)
  \item \textbf{(a)} Both are upper bounds and need not be tight: the first
        algorithm might be $\bigTh{n}$ and merely \emph{described} as
        $\bigO{n^{3}}$. \textbf{(b)} Even if both bounds are tight, the
        constants and the range of $n$ decide which is faster in practice ---
        $0.001n^{3}$ beats $1000n^{2}$ for every $n$ below $10^{6}$.
\end{ans}

%---------------------------------------------------------------------
\ansch{4 --- Correctness: Loop Invariants}
\begin{ans}
  \item Invariant: \emph{at the start of the iteration for $i$,
        \textit{total} $= \sum_{k=1}^{i-1} A[k]$.} Initialisation:
        $i = 1$ and the empty sum is 0. Maintenance: adding $A[i]$ extends the
        sum by one term. Termination: $i = n+1$ gives
        $\sum_{k=1}^{n}A[k]$. The measure is $n + 1 - i$, a non-negative
        integer strictly decreasing each iteration.
  \item A program that overwrites $A[1 \mathinner{\ldotp\ldotp} k]$ with
        zeros satisfies it: the prefix is sorted at every step. Repair:
        \emph{$A[1\mathinner{\ldotp\ldotp}k]$ is sorted \textbf{and} contains
        the $k$ smallest elements of the original array, and
        $A[k+1\mathinner{\ldotp\ldotp}n]$ holds the rest.} The second clause is
        what forbids inventing values.
  \item Initialisation and maintenance give \emph{partial} correctness --- if
        the loop halts, the invariant plus the exit condition give the
        postcondition. They say nothing about halting. What is missing is a
        \emph{loop variant}: an integer-valued measure that strictly decreases
        and is bounded below. Section~4.6's binary search has a valid invariant
        and no such measure, and it ran for ever on 13{,}228 of 20{,}000 trials.
\end{ans}

%---------------------------------------------------------------------
\ansch{5 --- Recurrences I}
\begin{ans}
  \item The tree is a \emph{path}: one node at each level, doing $n$, $n-1$,
        $n-2$, \ldots\ work, with $n$ levels. The total is
        $\sum_{k=1}^{n} k = n(n+1)/2 = \bigTh{n^{2}}$. It is not
        $\bigTh{n\lg n}$ because the argument shrinks by a constant
        \emph{amount}, not a constant \emph{fraction} --- so the depth is $n$
        rather than $\lg n$.
  \item Guess $T(n) \le cn^{2}$. Then
        $T(n) \le 2c(n/2)^{2} + n^{2} = cn^{2}/2 + n^{2}$, which is
        $\le cn^{2}$ provided $n^{2} \le cn^{2}/2$, i.e.\ $\mathbf{c \ge 2}$.
        With a base case at $n_{0} = 1$ and $c = \max\{2, T(1)\}$ the induction
        closes. (This is case 3 of the master theorem, confirming the answer.)
  \item A 1-to-99 split divides the problem by the constant factors 100 and
        $100/99$, so the depth is $\log_{100/99} n = \bigTh{\lg n}$ and every
        full level does $n$ work: $\bigTh{n\lg n}$, with a larger constant. A
        1-to-$(n-1)$ split removes one element per level, so the depth is $n$
        and the total is $\bigTh{n^{2}}$. \textbf{Constant fraction, not
        constant amount} is the distinction.
\end{ans}

%---------------------------------------------------------------------
\ansch{6 --- Recurrences II: The Master Theorem}
\begin{ans}
  \item $a = 4$, $b = 2$, $f(n) = n$. Watershed $n^{\log_{2}4} = n^{2}$. Since
        $f(n) = n = \bigO{n^{2-\varepsilon}}$ with $\varepsilon = 1$, this is
        \textbf{case 1}, and $T(n) = \bigTh{n^{2}}$. The leaves dominate; the
        combining work is irrelevant.
  \item \textbf{None.} Watershed $n^{\log_{2}2} = n$, and $f(n) = n/\lg n$ is
        \emph{smaller} than $n$ --- but only by a logarithmic factor, and case 1
        needs $f = \bigO{n^{1-\varepsilon}}$ for a \emph{constant}
        $\varepsilon > 0$. Since $n/\lg n = \omega(n^{1-\varepsilon})$ for
        every such $\varepsilon$, no case applies. (A recursion tree gives
        $\bigTh{n \lg\lg n}$.)
  \item In case 2 every level of the tree does the same total work, so the
        total is (work per level) $\times$ (number of levels), and the number
        of levels is $\lg_{b} n = \bigTh{\lg n}$. In case 1 the work grows
        geometrically towards the leaves and the last level dominates; in
        case 3 it shrinks geometrically and the root dominates. Only a tie
        makes every level count, and counting levels is where a logarithm comes
        from.
\end{ans}

%---------------------------------------------------------------------
\ansch{7 --- Brute Force}
\begin{ans}
  \item The sum is $\sum_{p=1}^{n} p(p+1)/2$. For $n = 100$ this is
        $\tfrac12\bigl(\tfrac{100\cdot101\cdot201}{6} + \tfrac{100\cdot101}{2}\bigr)
        = \tfrac12(338{,}350 + 5{,}050) = \mathbf{171{,}700}$. Against
        $n^{3}/6 = 166{,}667$ that is 3\% high, as it should be --- the
        leading term is $n^{3}/6$ and the next is $n^{2}/2$.
  \item The $n^{3} \to n^{2}$ step removed \emph{redundant arithmetic inside}
        the search: the same set of $\binom{n}{2}$ candidate blocks is still
        examined, each now in $\bigO{1}$ instead of $\bigO{n}$. The
        $n^{2} \to n$ step changed \emph{the search itself} --- Kadane's
        observation about blocks ending at $j$ means pairs are never enumerated
        at all. The first is careful programming; the second needs an idea
        about the problem.
  \item Ship the cubic version when $n$ is \emph{bounded by something real} and
        small --- a board of at most 64 squares, a form with at most 50 fields
        --- and the routine is off the critical path. Measured, it costs 79
        $\mu$s at $n = 100$. \textbf{The fact that would change my mind:} that
        $n$ is user-supplied or unbounded, because then the bound is a
        prediction rather than a fact.
\end{ans}

%---------------------------------------------------------------------
\ansch{8 --- Divide and Conquer}
\begin{ans}
  \item Because merge sort's control flow does not depend on the data. The
        split is always into halves and \algname{Merge} always performs
        $\bigTh{n_{1}+n_{2}}$ work, whatever the values. Insertion sort's inner
        loop runs a number of times that depends on the number of inversions,
        which is 0, $\approx n^{2}/4$ or $\approx n^{2}/2$ depending on the
        input --- so it has three analyses and merge sort has one.
  \item The minimum $\min(n_{1},n_{2})$ is attained when one run is entirely
        smaller than the other, as on \textbf{already-sorted} input, where the
        loop exhausts the left run first. The maximum $n_{1}+n_{2}-1$ is
        attained when the runs \textbf{interleave}, as on random input, so that
        both are nearly exhausted together. Measured: $0.496\,n\lg n$ against
        $0.940\,n\lg n$ --- almost exactly the predicted factor of 2.
  \item They measure different things. \chref{role}'s 65 is where insertion
        sort beats merge sort on a \emph{whole array}, timed from cold. The
        cutoff of 32 is where insertion sort beats \emph{continued recursion}
        on a subarray that is already in cache and already part of a larger
        sort. The second has less overhead to beat, so it turns over sooner.
        Either is a plausible answer; the point is that they are not the same
        quantity.
\end{ans}

%---------------------------------------------------------------------
\ansch{9 --- Quicksort}
\begin{ans}
  \item On sorted input the pivot (the last element) is the largest, so
        \algname{Partition} compares it with all $k-1$ others and splits
        $(k-1, 0)$. The subproblem sizes are therefore $n-1, n-2, \ldots, 1$,
        and the comparisons are $(n-1) + (n-2) + \cdots + 1 = n(n-1)/2$.
        Measured $0.9990 \to 0.9999$ of $n^{2}/2$.
  \item ``Average case'' averages over \emph{inputs} and so needs an assumption
        about the data; an adversary who knows your algorithm simply supplies
        an input outside that distribution --- and sorted data arrives by
        accident anyway. ``Expected'' averages over the algorithm's \emph{own
        coin flips} with the input held fixed, so it holds for every input
        including the adversary's. He would need to predict your random numbers,
        not your data.
  \item $1\,\text{MB} / 64\,\text{bytes} = 16{,}384$ frames, and the depth is
        $n-1$, so it crashes for $n > 16{,}384$ --- confirmed at
        $n = 32{,}000$. Doubling the stack buys $n = 32{,}768$: one more
        doubling of $n$, for twice the memory. It does not fix anything,
        because the depth grows linearly and the stack can only grow by
        constants. Recursing on the smaller side first bounds the depth at
        $\bigO{\lg n}$ and does fix it.
\end{ans}

%---------------------------------------------------------------------
\ansch{10 --- The Comparison Lower Bound}
\begin{ans}
  \item At least $4! = 24$ leaves, so height $\ge \lceil \lg 24 \rceil =
        \lceil 4.585 \rceil = \mathbf{5}$. And $\lg(4!) = 4.585$, so the bound
        says at least 5 comparisons in the worst case. (Five is in fact
        achievable, so the bound is tight here.)
  \item The step that fails is \emph{``a binary tree of height $h$ has at most
        $2^{h}$ leaves''} --- or rather, the premise that counting sort's
        execution is modelled by a binary decision tree at all. Counting sort
        uses a key as an \emph{array index}, an operation with $K$ possible
        outcomes rather than 2, so it extracts more than one bit per step and
        the tree is not binary. The theorem's hypothesis, not its conclusion,
        is what it escapes.
  \item The comparison count cannot improve by more than 1\%, so essentially
        none of \texttt{std::sort}'s 39\% comes from doing fewer comparisons.
        It comes from everything the comparison model does not charge for:
        sorting in place and so touching half as much memory, better locality,
        and fewer branch mispredictions. The lower bound tells you where the
        remaining performance is \emph{not}.
\end{ans}

%---------------------------------------------------------------------
\ansch{11 --- Selection and Order Statistics}
\begin{ans}
  \item Quicksort: $a = 2$, $b = 2$, watershed $n^{\log_{2}2} = n = f(n)$, so
        \textbf{case 2}, giving $\bigTh{n\lg n}$. Quickselect: $a = 1$,
        $b = 2$, watershed $n^{\log_{2}1} = n^{0} = 1$, far below $f(n) = n$,
        so \textbf{case 3}, giving $\bigTh{n}$. The $\lg n$ was the number of
        \emph{levels}; discarding one branch collapses the tree to a path whose
        work halves each step, and $n + n/2 + n/4 + \cdots < 2n$.
  \item The last-element pivot on sorted input is always the largest, so the
        range shrinks from the top by one each iteration. To reach the median
        the range must fall from $n$ to $n/2$, costing
        $n + (n-1) + \cdots + n/2 = \frac{(n + n/2)}{2}\cdot\frac{n}{2}
        = 3n^{2}/8$. Measured $0.7500 \times n^{2}/2$, which is the same
        number.
  \item With groups of 5, the pivot exceeds at least $3 \cdot \tfrac12 \cdot
        \tfrac{n}{5} = \tfrac{3n}{10}$ elements, so the recursion is on at most
        $7n/10$ and the recurrence is $T(n/5) + T(7n/10) + \bigTh{n}$ with
        $\tfrac15 + \tfrac{7}{10} = \tfrac{9}{10} < 1$. With groups of 3 the
        pivot exceeds only $2\cdot\tfrac12\cdot\tfrac{n}{3} = \tfrac{n}{3}$, so
        the recurrence is $T(n/3) + T(2n/3) + \bigTh{n}$ and
        $\tfrac13 + \tfrac23 = 1$. The substitution $T(n) \le cn$ then gives
        $cn + an \le cn$, which is false for every $c$, and the true bound is
        $\bigTh{n\lg n}$.
\end{ans}

%---------------------------------------------------------------------
\ansch{12 --- Heaps and Search Trees}
\begin{ans}
  \item Because the shape is enforced by the \emph{addressing scheme}, not
        maintained by an algorithm. A heap is stored in an array with node $i$'s
        children at $2i+1$ and $2i+2$, and a new element goes at index $n$.
        There is no insertion order that produces an unbalanced heap, because
        there is no array layout in which an unbalanced heap could be
        represented. A BST has no such constraint, which is why it needs
        rotations and a heap does not.
  \item The loose step is charging \emph{every} node the cost of the
        \emph{worst} node. In fact at most $\lceil n/2^{j+1}\rceil$ nodes have
        height $j$, so the total is
        $\bigO{n\sum_{j\ge0} j/2^{j}}$. And
        $\sum_{j \ge 0} j/2^{j} = \mathbf{2}$, from
        $\sum_{j\ge0} jx^{j} = x/(1-x)^{2}$ at $x = 1/2$. So the true total is
        $\bigO{2n} = \bigTh{n}$.
  \item \textbf{``In what order were the keys inserted?''} If they were
        shuffled, the height is about $2.3\lg n$ and the benchmark is
        meaningless for sorted input --- which is what real data often is.
        Measured, the same keys in sorted order gave height exactly $n$ and
        searches 650$\times$ slower.
\end{ans}

%---------------------------------------------------------------------
\ansch{13 --- Hashing}
\begin{ans}
  \item \emph{Each key is equally likely to hash to any of the $m$ slots,
        independently of the others.} It is a claim about the hash \emph{and}
        the keys because the hash function is deterministic: given the keys,
        the slot of each is fixed, and whether they scatter depends entirely on
        the relationship between the function and the key set. Section~13.3's
        $k \bmod 2^{20}$ is a perfectly good hash for unstructured keys and
        uses 1.6\% of the table for keys with a stride of 64.
  \item Searching for $x_{i}$, the elements examined first are those inserted
        after it that hashed to the same slot; for $j > i$ that happens with
        probability $1/m$. Averaging over $i$:
        $1 + \frac{1}{nm}\sum_{i=1}^{n}(n-i) = 1 + \frac{\alpha}{2} -
        \frac{\alpha}{2n}$. At $\alpha = 4$ this is $1 + 2 = 3$ for large $n$;
        measured $\mathbf{2.998}$.
  \item The \textbf{distribution} of chain lengths --- specifically the longest
        chain and the mean probe count --- not the mean chain length. The mean
        chain length is $n/m$ \emph{by definition}, whatever the hash does,
        because every key lands somewhere. Only the spread distinguishes a hash
        that scatters from one that piles up: measured, longest chain 8 against
        62, and mean probes 1.5 against 31.0.
\end{ans}

%---------------------------------------------------------------------
\ansch{14 --- Greedy Algorithms}
\begin{ans}
  \item \emph{Greedy-choice property:} some optimal solution contains the
        locally best first choice. \emph{Optimal substructure:} an optimal
        solution contains optimal solutions to the subproblems remaining after
        that choice. \textbf{A problem with the second and not the first:} 0/1
        knapsack. Its optimal solution does contain optimal sub-solutions, so
        dynamic programming works (\chref{dptwo}); but taking the highest
        value-per-weight item first can be wrong, so greedy fails. Coin change
        with $\{1,3,4\}$ is the same story.
  \item \textbf{8.} Greedy takes $4+3+1$ = three coins; optimal is $4+4$ = two.
        (The first failure is 6, as the chapter says; 7 is $4+3$ for both.)
        The counter-examples are small because greedy's mistake is made at the
        \emph{first} step, so the smallest instance exhibiting it is the
        smallest amount where the largest coin is the wrong first choice ---
        and that is bounded by the largest denomination.
  \item Huffman assigns each symbol a \emph{whole number} of bits, with a
        minimum of 1. A symbol of probability 0.95 carries only
        $-\lg 0.95 = 0.074$ bits of information, and Huffman must spend 1 on
        it. The algorithm is provably optimal \emph{among prefix codes}; the
        restriction to integer code lengths is what costs the 0.71 bits, and
        arithmetic coding, which can spend fractional bits, removes it.
\end{ans}

%---------------------------------------------------------------------
\ansch{15 --- Dynamic Programming I}
\begin{ans}
  \item Optimal substructure and overlapping subproblems. \textbf{A problem
        with the first and not the second:} merge sort --- sorting optimally
        requires sorting the halves, but the two halves are different
        subproblems and no subproblem ever recurs, so memoising it buys exactly
        nothing. Binary search is another.
  \item $T(0) = 1$ and $T(n) = 1 + \sum_{j=0}^{n-1} T(j)$. By induction,
        assuming $T(j) = 2^{j}$ for all $j < n$:
        $T(n) = 1 + \sum_{j=0}^{n-1}2^{j} = 1 + (2^{n}-1) = 2^{n}$. At
        $n = 22$ this is $4{,}194{,}304$, which is what the program counted ---
        exactly, not approximately.
  \item $\bigTh{mnk}$: (number of distinct subproblems) $\times$ (work per
        subproblem). For rod cutting there are $n$ subproblems and the maximum
        is over at most $n$ choices, giving $\bigTh{n \cdot n} =
        \bigTh{n^{2}}$. Note the measured caveat: with a price table capped at
        30 entries the inner loop runs at most 30 times, so the \emph{measured}
        behaviour was $\bigTh{30n}$ --- linear. Check the parameters before
        quoting the bound.
\end{ans}

%---------------------------------------------------------------------
\ansch{16 --- Dynamic Programming II}
\begin{ans}
  \item $c[i][j] = c[i-1][j-1]+1$ when $x_{i} = y_{j}$, else
        $\max(c[i-1][j], c[i][j-1])$, with $c[i][0] = c[0][j] = 0$. No maximum
        is needed in the equal case because including the matching character
        can never hurt: given any common subsequence of $X_{1..i}, Y_{1..j}$ not
        ending in it, appending it yields one at least as long. So some LCS
        does end in that character, and the rest is an LCS of the two shortened
        prefixes.
  \item A running time is polynomial when polynomial in the \emph{size} of the
        input, and a number's size is its bit length (\chref{cases}). $W$ is
        one number of $b \approx \lg W$ bits, so $\bigTh{nW} =
        \bigTh{n2^{b}}$ --- exponential in $b$. Measured: eight doublings of
        $W$, each adding one bit, gave time ratios of mean $2.02$. This is the
        same phenomenon as \chref{cases}'s trial division, whose eight ratios
        averaged $2.02$ per two bits, reached from the opposite direction.
  \item \textbf{Appropriate:} computing only the \emph{length} of an LCS, or
        only the edit \emph{distance}, where the answer is the final table
        entry. \textbf{Not appropriate:} recovering the subsequence or the
        alignment itself, which requires walking backwards through the whole
        table. What distinguishes them is whether the answer is a single value
        or a \emph{path}; a path needs the table that produced it. Hirschberg's
        algorithm gets both, at twice the code.
\end{ans}

%---------------------------------------------------------------------
\ansch{17 --- Amortised Analysis}
\begin{ans}
  \item Copying happens at lengths $1, 2, 4, \ldots$, the powers of two up to
        $n$, and the copy at length $2^{k}$ moves $2^{k}$ elements. So the
        total is $\sum_{k=0}^{\lfloor\lg n\rfloor} 2^{k} =
        2^{\lfloor\lg n\rfloor+1}-1 < 2n$. The measured 1.024 is just over
        $1n$, not $2n$, because when $n$ is a power of two the sum is exactly
        $n-1$; the bound of $2n$ is the worst position within a doubling
        interval.
  \item \textbf{Cheap:} $c = 1$, \textit{len} rises by 1 so $\Phi$ rises by 2,
        giving $\hat c = 1 + 2 = 3$. \textbf{Expensive} (resize from $m$ to
        $2m$ with \textit{len} $= m$): $c = m+1$; before,
        $\Phi = 2m - m = m$; after, $\Phi = 2(m{+}1) - 2m = 2$; so
        $\hat c = (m+1) + 2 - m = 3$. They \emph{must} come out equal, or the
        potential function would not be doing its job --- the whole point is to
        choose $\Phi$ so that the expensive operation's real cost is exactly
        paid for by the potential the cheap ones accumulated.
  \item Any \textbf{latency-bounded} context: a 16 ms frame deadline, an
        interrupt handler, an audio callback, a trading system. Measured, the
        amortised cost was 1.049 copies per append and the worst single append
        copied \textbf{524{,}288} elements --- half a million times the
        average. And only 20 of a million appends were expensive, so a
        latency histogram with a thousand samples would very likely contain
        none of them, and the system would look fine until it was not.
\end{ans}

%---------------------------------------------------------------------
\ansch{18 --- Graphs}
\begin{ans}
  \item Each edge $\{u,v\}$ contributes 1 to $\deg(u)$ and 1 to $\deg(v)$, so
        summing degrees counts every edge twice. Theorem~18.1 uses it to turn
        the per-vertex inner loop into a bound in $|E|$: the BFS inner loop
        runs $\deg(u)$ times for each dequeued $u$, and
        $\sum_{u}\deg(u) = 2|E|$ gives the $\bigTh{|V|+|E|}$.
  \item $|E| = 3 \times 10^{6}$. \textbf{Matrix:} $|V|^{2} = 10^{12}$ bytes
        $= 1$ TB --- impossible. \textbf{List:} about
        $(10^{6} \cdot 8) + (6 \times 10^{6} \cdot 4) \approx 32$ MB. BFS is
        $\bigTh{|V|^{2}} = 10^{12}$ steps from the matrix and
        $\bigTh{|V|+|E|} = 7 \times 10^{6}$ from the list --- a factor of
        140{,}000, and the matrix cannot be built anyway.
  \item That a crossover density quoted without an implementation is
        meaningless. The same two \emph{data structures} gave a crossover at
        0.26 with one implementation of the matrix and no crossover at all with
        another, and the gap between a good and a poor implementation of the
        matrix (8--20$\times$) was larger than the gap between the two
        structures. Quote the crossover for \emph{your} code, measured.
\end{ans}

%---------------------------------------------------------------------
\ansch{19 --- Minimum Spanning Trees}
\begin{ans}
  \item \emph{Let $A$ be a subset of some MST and let $(S, V\setminus S)$ be a
        cut that no edge of $A$ crosses. Then any light edge for that cut is
        safe.} The hypothesis that makes the exchange safe is \textbf{that no
        edge of $A$ crosses the cut}: it guarantees that the edge $e'$ removed
        from the cycle --- which does cross the cut --- is not in $A$, so
        $A$ survives the exchange and remains inside the new MST.
  \item \textbf{Prim's cut} is (vertices in the tree so far) against the rest,
        so the light edge is the cheapest edge leaving the tree, and the
        structure needed is a \textbf{priority queue}. \textbf{Kruskal's cut}
        is between the two components that the next edge would join, so what is
        needed is a test for ``are these already connected?'' ---
        \textbf{disjoint-set union}. One theorem, two cuts, two data
        structures.
  \item It is wrong because $\alpha(n)$ is not constant --- it grows, and the
        measurement shows it growing, from 1.62 to 1.88. It does not matter
        because $\alpha$ grows so slowly that $\alpha(n) \le 4$ for every $n$
        that can be written down; reaching 5 needs an $n$ far larger than the
        number of atoms in the observable universe. It is also an
        \emph{amortised} bound, so a single \texttt{find} can still cost
        $\bigTh{\lg|V|}$.
\end{ans}

%---------------------------------------------------------------------
\ansch{20 --- Shortest Paths}
\begin{ans}
  \item The inequality $\delta(s,y) \le \delta(s,u)$, where $y$ lies on a
        shortest path to $u$. It says extending a path cannot make it cheaper,
        which requires non-negative weights. Without it, the remaining portion
        from $y$ to $u$ may have negative total weight, so a vertex finalised
        early can later prove reachable more cheaply --- and Dijkstra never
        revisits it. Measured in Section~20.4: $d(3) = 4$ where the answer
        is 3.
  \item Because a shortest path in a graph with no negative cycle uses at most
        $|V|-1$ edges (more would repeat a vertex, and the cycle so formed has
        non-negative weight and can be removed). Each pass propagates correct
        distances one more edge along every path, so $|V|-1$ passes suffice.
        \textbf{The $|V|$-th pass detects a negative cycle}: if any distance
        still improves, some path is using $|V|$ or more edges and still
        getting shorter, which only a negative cycle permits.
  \item That a worst-case bound describes the worst case and not the typical
        one, and that both facts matter. Bellman--Ford's $\bigO{|V||E|}$ is
        real --- the path graph attained it, 789$\times$ slower than Dijkstra
        --- and on random graphs the early exit made it indistinguishable. A
        benchmark on random inputs alone would have reported them equal; the
        bound is what told us to go looking for the input where they are not.
\end{ans}

%---------------------------------------------------------------------
\ansch{21 --- String Matching}
\begin{ans}
  \item For $P = \texttt{aabaaab}$:
        $\pi = [0, 1, 0, 1, 2, 2, 3]$. (At $i=5$, \texttt{aa} is a border of
        \texttt{aabaa}; at $i=7$, \texttt{aab} is a border of
        \texttt{aabaaab}.) After a mismatch having matched 6 characters,
        $k \gets \pi[6] = 2$ --- the pattern shifts so that its first 2
        characters align with the last 2 matched, a shift of $6-2 = 4$
        positions.
  \item Because the inner loop only ever \emph{decreases} $k$, and $k$ is
        increased by at most 1 per outer iteration. Take $\Phi = k$: it starts
        at 0, never goes negative, and rises by at most $n$ in total across the
        run, so it can fall at most $n$ times in total. The inner loop
        therefore executes at most $n$ times \emph{over the whole run}, not per
        iteration --- the same amortised argument as \algname{Build-Heap} and
        the dynamic array.
  \item \textbf{Naive:} every alignment matches the $m-1$ leading
        \texttt{a}'s and fails on the final \texttt{b}, so each of the $n$
        shifts costs $m$ comparisons --- the worst case, attained exactly.
        \textbf{Rabin--Karp:} it never compares characters at all unless the
        hashes agree. Every window of $\texttt{a}^{n}$ has the same hash, and
        $\texttt{a}^{m-1}\texttt{b}$ has a different one, so every window is
        rejected in one integer comparison. The two algorithms fail on
        different things: naive on \emph{long partial matches}, Rabin--Karp on
        \emph{hash collisions}.
\end{ans}

%---------------------------------------------------------------------
\ansch{22 --- Complexity Classes}
\begin{ans}
  \item \textbf{P} is the class of decision problems solvable in polynomial
        time. \textbf{NP} is the class whose yes-instances have a certificate
        verifiable in polynomial time. $\mathrm{P} \subseteq \mathrm{NP}$
        because if you can \emph{solve} a problem in polynomial time you can
        verify any proposed certificate in polynomial time by ignoring it and
        solving the instance yourself.
  \item From a \textbf{known hard problem to $X$}: $\text{SAT} \le_{P} X$.
        That shows $X$ is at least as hard as SAT, since a fast algorithm for
        $X$ would give one for SAT. The other direction, $X \le_{P}
        \text{SAT}$, proves only that $X$ is no \emph{harder} than SAT --- which
        is true of every problem in NP and says nothing about $X$'s difficulty.
  \item $\bigTh{nW}$ is polynomial in the \emph{value} $W$ and exponential in
        its \emph{length} $\lg W$, so it is pseudo-polynomial, not polynomial,
        and does not place knapsack in P. \chref{dptwo} measured this: with $n$
        fixed at 200, eight doublings of $W$ --- each adding exactly one bit of
        input --- multiplied the running time by a mean factor of 2.02. An
        algorithm whose time doubles when the input grows by one bit is
        exponential in the input size.
\end{ans}

%---------------------------------------------------------------------
\ansch{23 --- The Semester Project}
\begin{ans}
  \item For example: \emph{``The best insertion-sort cutoff for 64-byte records
        is larger than the best cutoff for 4-byte integers, on the same
        machine.''} It predicts a direction, it names the quantity to measure,
        and the result ``it is smaller'' or ``there is no difference'' would
        refute it. Compare with ``I will compare sorting algorithms'', which no
        measurement could contradict.
  \item \textbf{Your noise level and your clock resolution.} If repeated runs of
        the \emph{same} measurement vary by $\pm 15\%$, then 3.6--4.4 is
        consistent with 4 and the answer is yes. If they vary by $\pm 1\%$,
        a spread of 3.6 to 4.4 is a real effect and something else is going on.
        Also check that the smallest time is well above the clock resolution.
        A ratio without a noise estimate is not evidence either way.
  \item For example: \textbf{(a)} uniformly random keys, probing the ideal
        case the analysis assumes; \textbf{(b)} keys with a fixed stride, such
        as multiples of 64, probing whether the hash uses the whole key ---
        \chref{hashing} measured 31 probes instead of 1.5 here; \textbf{(c)}
        keys deliberately chosen to collide under your hash, probing the
        adversarial worst case. Several other answers are good: a realistic
        sample of your own data is often the most informative of all.
\end{ans}
